\documentclass[10pt,a4paper]{amsart}
\usepackage[margin=19mm]{geometry}
\usepackage{amsmath,amssymb,mathtools}
\usepackage[scr=rsfs,cal=euler]{mathalfa}
\usepackage{xfrac}
\usepackage[unicode,hidelinks,bookmarks=false]{hyperref}
\newcommand{\ie}{i.e.\ }
\newcommand{\cf}{cf.\ }
\newcommand{\resp}{respectively\ }
\newcommand{\Subsection}{\subsection}
\providecommand{\nobreakdash}{\nobreak\hbox{-}}
\setlength{\emergencystretch}{2em}
\multlinegap0pt
\usepackage{bm}
\usepackage{amssymb}
\usepackage{mathtools}
\usepackage{tikz}
\newtheorem{lemma}{Lemma}
\title{The Ergodicity of Geodesic Flows on Rank One Manifolds of Nonpositive Curvature}
\author{Fei Liu, Xiaokai Liu}
\date{Source: arXiv:2609.12381v1 (CC BY 4.0)}
\begin{document}
\maketitle

\noindent\emph{Source and licence.} The Ergodicity of Geodesic Flows on Rank One Manifolds of Nonpositive Curvature. Version arXiv:2609.12381v1 (CC BY 4.0), distributed by arXiv under the Creative Commons Attribution 4.0 International licence, http://creativecommons.org/licenses/by/4.0/, as recorded in the arXiv licence metadata for this exact version and shown on the abstract page https://arxiv.org/abs/2609.12381v1. Full licence notice: \texttt{licenses/arxiv-2609.12381-CC-BY-4.0.txt}.\par\smallskip

\noindent\textit{Editorial scope.} Counted unit: the article's lemma lem\_3\_2 (source0.tex lines 195--202). The article's complete original proof of it is delivered with it (source0.tex lines 204--284, one \texttt{proof} environment of 81 lines). No mathematics is written, repaired or re-derived: the statement and the proof are the article's own lines, in the article's own notation, and no retained line is substituted or re-flowed.  No further source-local context is needed: every cross-reference of every retained line resolves inside the two delivered spans.  Nothing was omitted from any retained span, so the package carries no omission record.  No retained line contains a citation, so the wrapper carries no bibliography and no bibliography entry.  No retained line contains a figure environment, an image-inclusion command or drawing code, so no image is delivered and the package declares no asset dependency.  Editorial formatting only, in addition: an AMS \texttt{amsart} preamble that restates the article's own declarations and macros used by the retained lines, namely the environments \texttt{lemma} on separate counters, together with the standard packages those lines need.  The retained environments print in this document as Lemma 1 in order of appearance, whereas the article prints the same items with the numbering of its own full document.  The complete native source of the article as distributed in the arXiv e-print for this exact version is bundled unchanged as \texttt{sources/210110/source0.tex}.
\begin{lemma}\label{lem_3_2}
	For each $x\in M$, let
	\begin{displaymath}
		\Lambda_{x} = \left\{\bm{v}\in T^{1}_{x}M \mid K(c_{\bm{v}}(t))=0,~\forall t\in\mathbb{R}\right\}.
	\end{displaymath}

	If $x\notin F_{\infty}$, then $\Lambda_x$ is a finite set, thus has zero angular measure on $T_x^1M\simeq S^1$.
\end{lemma}

\begin{proof}
	Take any point $x\notin F_{\infty}$. If $K(x)<0$, $\Lambda_{x}=\emptyset$ by definition. When $K(x)=0$, $0$ is a local maxima, $\nabla K(x)=0$. Since $x\notin F_{\infty}$, we can find the minimal $m\ge 2$, such that $\nabla^{m}K(x)\neq 0$.

	By definition, for any $\bm{v}\in \Lambda_{x}$, $K(c_{\bm{v}}(t))=0$ for any $t\in\mathbb{R}$. Thus,
	\begin{displaymath}
		\dfrac{\mathrm{d}^n}{\mathrm{d}t^{n}} K(c_{\bm{v}}(t))=0,\quad\forall n\in\mathbb{N}.
	\end{displaymath}

	In particular, $\dfrac{\mathrm{d}^{m}}{\mathrm{d}t^{m}} K(c_{\bm{v}}(t))=0$.

	Now claim that
	\begin{displaymath}
		\nabla^{n}K(x)(\overbrace{\bm{v},\ldots,\bm{v}}^{n})=\left.\dfrac{\mathrm{d}^{n}}{\mathrm{d}t^{n}}\right\rvert_{t=0}K(c_{\bm{v}}(t)).
	\end{displaymath}

	Prove this claim by induction. When $n=1$, by chain rule
	\begin{displaymath}
		\left.\dfrac{\mathrm{d}}{\mathrm{d}t}\right\rvert_{t=0}K(c_{\bm{v}}(t))=\langle \nabla K(c_{\bm{v}}(t)),c'_{\bm{v}}(t)\rangle\rvert_{t=0}=\langle \nabla K(x),\bm{v}\rangle=\nabla K(x)(\bm{v}).
	\end{displaymath}

	Suppose the claim is true for $n=k$, then we can compute that
	\begin{displaymath}
		\begin{aligned}
			\dfrac{\mathrm{d}^{k+1}}{\mathrm{d}t^{k+1}}K(c_{\bm{v}}(t)) & = \dfrac{\mathrm{d}}{\mathrm{d}t}(\nabla^{k} K(c_{\bm{v}}))(t)(\overbrace{c'_{\bm{v}}(t),\ldots,c'_{\bm{v}}(t)}^{k})                         \\
			                                                            & =(\nabla_{c'_{\bm{v}}(t)}\nabla^{k}K(c_{\bm{v}}(t)))(\overbrace{c'_{\bm{v}}(t),\ldots,c'_{\bm{v}}(t)}^{k})                                   \\
			                                                            & +\sum_{j=1}^k \nabla^{k}K(c_{\bm{v}}(t))(c'_{\bm{v}}(t),\ldots,\overbrace{\nabla_{c'_{\bm{v}}(t)}c'_{\bm{v}}(t)}^{j},\ldots, c'_{\bm{v}}(t)) \\
		\end{aligned}
	\end{displaymath}

	Given $c_{\bm{v}}(t)$ is a geodesic, $\nabla_{c'_{\bm{v}}(t)}c'_{\bm{v}}(t)=0$, terms in the summation all vanish. Therefore, we obtain the required formula:
	\begin{displaymath}
		\begin{aligned}
			\left.\dfrac{\mathrm{d}^{k+1}}{\mathrm{d}t^{k+1}}\right\rvert_{t=0}K(c_{\bm{v}}(t)) & =\nabla_{c'_{\bm{v}}(t)}\nabla^{k}K(c_{\bm{v}}(t))(\overbrace{c'_{\bm{v}}(t),\ldots,c'_{\bm{v}}(t)}^{k})\rvert_{t=0} \\
			                                                                                    & =\nabla^{k+1}K(c_{\bm{v}}(t))(\overbrace{c'_{\bm{v}}(t),\ldots,c'_{\bm{v}}(t)}^{k+1})\rvert_{t=0}                    \\
			                                                                                    & =\nabla^{k+1} K(x)(\overbrace{\bm{v},\ldots,\bm{v}}^{k+1}).
		\end{aligned}
	\end{displaymath}

	Now let $P_{m,x}(\bm{u})\coloneq(\nabla^{m}K(x))(\bm{u},\ldots,\bm{u})$, we can compute the component of $P_{m,x}$ in local coordinate $(x_1,x_2)$ near the point $x$. Start from the first order, we have
	\begin{displaymath}
		\nabla K(x)=\dfrac{\partial K(x)}{\partial x_{1}}\mathrm{d}x_{1} + \dfrac{\partial K(x)}{\partial x_{2}}\mathrm{d}x_2=0.
	\end{displaymath}

	Now the second order covariant derivative $\nabla^{2}K(x)$ is given by
	\begin{displaymath}
		\nabla_{j}\nabla_{i}K(x)=\dfrac{\partial^2 K(x)}{\partial x_{j}\partial x_{i}}-\Gamma^{k}_{ij}\dfrac{\partial K}{\partial x_{k}}=\dfrac{\partial^2 K(x)}{\partial x_{j}\partial x_{i}}.
	\end{displaymath}

	By induction on $l<m$, we obtain that,
	\begin{displaymath}
		0={(\nabla^{l} K(x))}_{i_{1}\cdots i_{l}}=\dfrac{\partial^{l} K(x)}{\partial x_{i_1}\cdots\partial x_{i_l}},\qquad i_1,\ldots,i_l=1,2,
	\end{displaymath}
	provided that all lower order derivatives vanish at $x$. In particular, when $l=m$
	\begin{displaymath}
		{\left(\nabla^{m} K(x)\right)}_{i_{1} \cdots i_{m}}=\dfrac{\partial^{m} K(x)}{\partial x_{i_{1}}\cdots\partial x_{i_{m}}}+\sum_{0\leq l\leq m-1} A\cdot B,\qquad i_1,\ldots,i_m=1,2.
	\end{displaymath}

	In the summation, $A$ stands for lower orders of partial derivatives of $K(x)$, which all vanish at $x$, and $B$ stands for the Christoffel symbols and their derivatives. Therefore, we get
	\begin{equation}\label{eq_3_1}
		{(\nabla^{m} K(x))}_{i_1\cdots i_m}=\dfrac{\partial^m K(x)}{\partial x_{i_1}\cdots\partial x_{i_m}}.
	\end{equation}

	Choose $\bm{u}\in T^{1}_x M$. In the local geodesic normal coordinate with the orthonormal frame $\{\bm{e}_{1}(y),\bm{e}_{2}(y)\}$ for $y$ near $x$, set $\bm{u}=u_{1}\,\bm{e}_{1}(x)+u_{2}\,\bm{e}_{2}(x)=u_{1}\,\partial_{x_1}+u_{2}\,\partial_{x_2}$, we can write
	\begin{equation}\label{eq_3_2}
		P_{m,x}(\bm{u})=\nabla^{m} K(x)(\overbrace{\bm{u},\ldots,\bm{u}}^{m})=\sum_{i_{1},\ldots,i_{m}=1}^{2}\dfrac{\partial^{m}K(x)}{\partial x_{i_1}\cdots\partial x_{i_m}}u_{i_1}\ldots u_{i_m}.
	\end{equation}

	Equation~\ref{eq_3_2} shows that $P_{m,x}(\bm{u})$ is a nonzero homogeneous polynomial of degree $m$ in $u_1$ and $u_2$.

	Now, for any $\bm{v}\in \Lambda_x$, by definition, the geodesic $c_{\bm{v}}$ is flat, hence
	\begin{displaymath}
		P_{m,x}(\bm{v})=\dfrac{\mathrm{d}^{m}}{\mathrm{d}t^{m}}K(C_{\bm{v}}(t))=0.
	\end{displaymath}

	On $T^1_x M\simeq S^1$, let $\bm{v}=\cos\theta\bm{e}_1(x)+\sin\theta\bm{e}_2(x)$. By homogeneity, the restriction of $P_{m,x}(\bm{v})$ to the unit circle now becomes a trigonometric polynomial of $\theta$ of degree $m$:

	\begin{equation}\label{eq_3_3}
		P_{m,x}(\bm{v}(x,\theta))=\sum_{i=0}^{m} a_{i}(x)\cos^{m-i}\theta\sin^{i}\theta,
	\end{equation}
	where $a_i(x)$ are smooth functions defined near $x$. For fixed $x$, there are only finitely many $\theta\in S^{1}$, such that $P_{m,x}=0$. Thus, $\Lambda_x$ is a finite set with $0$ angular measure.
\end{proof}

\end{document}
